% PRW 2024 Paper Template
% based on the CVPR template provided by Ming-Ming Cheng (https://github.com/MCG-NKU/CVPR_Template)
% modified and extended by Stefan Roth (stefan.roth@NOSPAMtu-darmstadt.de)
% modified and extended by Peter M.  Roth (peter.m.roth@vetmeduni.ac.at)

\documentclass[10pt,twocolumn,letterpaper]{article}

%%%%%%%%% PAPER TYPE  - PLEASE UPDATE FOR FINAL VERSION
\usepackage[review]{cvpr}      % To produce the REVIEW version
% \usepackage{cvpr}              % To produce the CAMERA-READY version
%\usepackage[pagenumbers]{cvpr} % To force page numbers, e.g. for an arXiv version

% Include other packages here, before hyperref.
\usepackage{graphicx}
\usepackage{amsmath}
\usepackage{amssymb}
\usepackage{booktabs}


% Support for easy cross-referencing
\usepackage[capitalize]{cleveref}
\crefname{section}{Sec.}{Secs.}
\Crefname{section}{Section}{Sections}
\Crefname{table}{Table}{Tables}
\crefname{table}{Tab.}{Tabs.}


%%%%%%%%% PAPER ID  - PLEASE UPDATE
\def\cvprPaperID{*****} % *** Enter the CVPR Paper ID here
\def\confName{PRW}
\def\confYear{2024}


\begin{document}

%%%%%%%%% TITLE - PLEASE UPDATE
\title{\LaTeX\ Author Guidelines for \confName~Proceedings}

\author{First Author\\
Institution1\\
Institution1 address\\
{\tt\small firstauthor@i1.org}
% For a paper whose authors are all at the same institution,
% omit the following lines up until the closing ``}''.
% Additional authors and addresses can be added with ``\and'',
% just like the second author.
% To save space, use either the email address or home page, not both
\and
Second Author\\
Institution2\\
First line of institution2 address\\
{\tt\small secondauthor@i2.org}
}
\maketitle


%%%%%%%%% ABSTRACT
\begin{abstract}
The ABSTRACT is to be in fully justified italicized text, at the top of the left-hand column, below the author and affiliation information. Use the word ``Abstract'' as the title, in 12-point Times, boldface type, centered relative to the column, initially capitalized. The abstract is to be in 10-point, single-spaced type. Leave two blank lines after the Abstract, then begin the main text. Look at previous \confName\ abstracts to get a feel for style and length.
\end{abstract}


%%%%%%%%% BODY TEXT
\section{Introduction}
\label{sec:intro}

This document gives a specification for the paper layout for submission to the Pattern Recognition Workshop 2023. The template is based on the the official CVPR 2022 template. This document is written in accordance to the specification and guidelines, and hence may be used as an example of how a paper should look like. 

\subsection{Language}

All manuscripts must be in English.


\subsection{Dual submission}

Please refer to the author guidelines on the \confName\ \confYear\ web page for a
discussion of the policy on dual submissions.


\subsection{Paper length} 

\confName\  papers may be of the following types (lengths includung references):

\begin{itemize}

\item Research paper (8 pages)
\item Overiew paper (8 pages)
\end{itemize}

Overlength papers will simply not be reviewed. This includes papers where the 
margins and formatting are deemed to have been significantly altered from those 
laid down by this style guide. 


% Teaser figure on first page
\begin{figure}[t]
  \centering
  \fbox{\rule{0pt}{1.8in} \rule{0.7\linewidth}{0pt}}
   %\includegraphics[width=0.8\linewidth]{egfigure.eps}

   \caption{Example of caption \& teaser figure on first page: It is set in 
   Roman so that mathematics (always set in Roman: $B \sin A = A \sin B$) may
   be included without an ugly clash.}
   \label{fig:onecol}
\end{figure}


\subsection{Review and Publication}

The reviewing process will be double-blind. Accepted papers will be presented at the workshop (talk or presentation) published in the workshop proceedings. 


\subsection{The ruler}
The \LaTeX\ style defines a printed ruler which should be present in the version submitted for review. The ruler is provided in order that reviewers may comment on particular lines in the paper without circumlocution. If you are preparing a document using a non-\LaTeX\ document preparation system, please arrange for an equivalent ruler to appear on the final output pages. The presence or absence of the ruler should not change the appearance of any other content on the page. The camera-ready copy should not contain a ruler. (\LaTeX\ users may use options of cvpr.sty to switch between different versions.)

\textbf{Reviewers}: Please note that the ruler measurements do not align well with lines in the paper -- this turns out to be very difficult to do well when the paper contains many figures and equations, and, when done, looks ugly. Just use fractional references (\eg, this line is $087.5$), although in most cases one would expect that the approximate location will be adequate. 


\subsection{Paper ID}
Make sure that the Paper ID from the submission system is visible in the version submitted for review (replacing the ``*****'' you see in this document). If you are using the \LaTeX\ template, \textbf{make sure to update paper ID in the appropriate place in the tex file}. 



% Example of short caption and using subfigures!
\begin{figure*}
  \centering
  \begin{subfigure}{0.68\linewidth}
    \fbox{\rule{0pt}{2in} \rule{.9\linewidth}{0pt}}
    \caption{An example of a subfigure.}
    \label{fig:short-a}
  \end{subfigure}
  \hfill
  \begin{subfigure}{0.28\linewidth}
    \fbox{\rule{0pt}{2in} \rule{.9\linewidth}{0pt}}
    \caption{Another example of a subfigure.}
    \label{fig:short-b}
  \end{subfigure}
  \caption{Example of a short caption, which should be centered.}
  \label{fig:short}
\end{figure*}


\section{Formatting your paper}
\label{sec:formatting}

All text must be in a two-column format. The total allowable size of the text area is $6\frac78$ inches (17.46 cm) wide by $8\frac78$ inches (22.54 cm) high. Columns are to be $3\frac14$ inches (8.25 cm) wide, with a $\frac{5}{16}$ inch (0.8 cm) space between them. The main title (on the first page) should begin 1 inch (2.54 cm) from the top edge of the page. The second and following pages should begin 1 inch (2.54 cm) from the top edge. On all pages, the bottom margin should be $1\frac{1}{8}$ inches (2.86 cm) from the bottom edge of the page for $8.5 \times 11$-inch paper; for A4 paper, approximately $1\frac{5}{8}$ inches (4.13 cm) from the bottom edge of the page.


\subsection{Margins and page numbering}

All printed material, including text, illustrations, and charts, must be kept within a print area $6\frac{7}{8}$ inches (17.46 cm) wide by $8\frac{7}{8}$ inches (22.54 cm) high. Page numbers should be in the footer, centered and $\frac{3}{4}$ inches from the bottom of the page. The review version should have page numbers, yet the final version submitted as camera ready should not show any page numbers. The \LaTeX\ template takes care of this when used properly.


\subsection{Type style and fonts}

Wherever Times is specified, Times Roman may also be used. If neither is available on your word processor, please use the font closest in appearance to Times to which you have access.

\textbf{MAIN TITLE}: Center the title $1\frac{3}{8}$ inches (3.49 cm) from the top edge of the first page. The title should be in Times 14-point, boldface type. Capitalize the first letter of nouns, pronouns, verbs, adjectives, and adverbs; do not capitalize articles, coordinate conjunctions, or prepositions (unless the title begins with such a word). Leave two blank lines after the title.

\textbf{AUTHOR NAME(s)} and \textbf{AFFILIATION(s)} are to be centered beneath the title and printed in Times 12-point, non-boldface type. This information is to be followed by two blank lines.

The \textbf{ABSTRACT} and \textbf{MAIN TEXT} are to be in a two-column format.

\textbf{MAIN TEXT}: Type main text in 10-point Times, single-spaced. Do NOT use double-spacing. All paragraphs should be indented 1 pica (approx.~$\frac{1}{6}$ inch or 0.422 cm). Make sure your text is fully justified---that is, flush left and flush right. Please do not place any additional blank lines between paragraphs.

Figure and table captions should be 9-point Roman type as in \cref{fig:onecol,fig:short}. Short captions should be centred.

\noindent Callouts should be 9-point Helvetica, non-boldface type. Initially capitalize only the first word of section titles and first-, second-, and third-order headings.

\textbf{FIRST-ORDER HEADINGS}: (For example, {\large \bf 1. Introduction}) should be Times 12-point boldface, initially capitalized, flush left, with one blank line before, and one blank line after.

\textbf{SECOND-ORDER HEADINGS}: (For example, { \bf 1.1. Database elements}) should be Times 11-point boldface, initially capitalized, flush left, with one blank line before, and one after. If you require a third-order heading (we discourage it), use 10-point Times, boldface, initially capitalized, flush left, preceded by one blank line, followed by a period and your text on the same line. 


\subsection{Footnotes}

Please use footnotes\footnote{This is what a footnote looks like. It often distracts the reader from the main flow of the argument.} sparingly. Indeed, try to avoid footnotes altogether and include necessary peripheral observations in the text (within parentheses, if you prefer, as in this sentence). If you wish to use a footnote, place it at the bottom of the column on the page on which it is referenced. Use Times 8-point type, single-spaced.


%-------------------------------------------------------------------------
\subsection{Cross-references}

For the benefit of author(s) and readers, please use the
{\small\begin{verbatim}
  \cref{...}
\end{verbatim}}  command for cross-referencing to figures, tables, equations, or sections.
This will automatically insert the appropriate label alongside the cross-reference as in this example:
\begin{quotation}
  To see how our method outperforms previous work, please see \cref{fig:onecol} and \cref{tab:example}.
  It is also possible to refer to multiple targets as once, \eg~to \cref{fig:onecol,fig:short-a}.
  You may also return to \cref{sec:formatting} or look at \cref{eq:also-important}.
\end{quotation}
If you do not wish to abbreviate the label, for example at the beginning of the sentence, you can use the
{\small\begin{verbatim}
  \Cref{...}
\end{verbatim}}
command. Here is an example:
\begin{quotation}
  \Cref{fig:onecol} is also quite important.
\end{quotation}


\subsection{References}

List and number all bibliographical references in 9-point Times, single-spaced, at the end of your paper. When referenced in the text, enclose the citation number in square brackets, for example~\cite{Authors14}. Where appropriate, include page numbers and the name(s) of editors of referenced books. When you cite multiple papers at once, please make sure that you cite them in numerical order like this \cite{Alpher02,Alpher03,Alpher05,Authors14b,Authors14}. If you use the template as advised, this will be taken care of automatically.

\begin{table}
  \centering
  \begin{tabular}{@{}lc@{}}
    \toprule
    Method & Frobnability \\
    \midrule
    Theirs & Frumpy \\
    Yours & Frobbly \\
    Ours & Makes one's heart Frob\\
    \bottomrule
  \end{tabular}
  \caption{Results.   Ours is better.}
  \label{tab:example}
\end{table}


\subsection{Illustrations, graphs, and photographs}

All graphics should be centered. In \LaTeX, avoid using the \texttt{center} environment for this purpose, as this adds potentially unwanted whitespace. Instead use
%
{\small\begin{verbatim}
  \centering
\end{verbatim}}
%
at the beginning of your figure. Please ensure that any point you wish to make is resolvable in a printed copy of the paper. Resize fonts in figures to match the font in the body text, and choose line widths that render effectively in print. Readers (and reviewers), even of an electronic copy, may choose to print your paper in order to read it. You cannot insist that they do otherwise, and therefore must not assume that they can zoom in to see tiny details on a graphic. When placing figures in \LaTeX, it's almost always best to use \verb+\includegraphics+, and to specify the figure width as a multiple of the line width as in the example below
%
{\small\begin{verbatim}
   \usepackage{graphicx} ...
   \includegraphics[width=0.8\linewidth]
                   {myfile.pdf}
\end{verbatim}
}


\subsection{Mathematics}

Please number all of your sections and displayed equations as in these examples:

\begin{equation}
  E = m\cdot c^2
  \label{eq:important}
\end{equation}

and

\begin{equation}
  v = a\cdot t.
  \label{eq:also-important}
\end{equation}

It is important for readers to be able to refer to any particular equation. Just because you did not refer to it in the text does not mean some future reader might not need to refer to it. It is cumbersome to have to use circumlocutions like ``the equation second from the top of page 3 column 1''. (Note that the ruler will not be present in the final copy, so is not an alternative to equation numbers).


\subsection{Miscellaneous}

The space after \eg, meaning ``for example'', should not be a sentence-ending space. So \eg is correct, {\em e.g.} is not. The provided \verb'\eg' macro takes care of this.

When citing a multi-author paper, you may save space by using ``et alia'', shortened to ``\etal'' (not ``{\em et.\ al.}'' as ``{\em et}'' is a complete word). If you use the \verb'\etal' macro provided, then you need not worry about double periods when used at the end of a sentence as in Alpher \etal. However, use it only when there are three or more authors.

Thus, the following is correct:
   ``Frobnication has been trendy lately.
   It was introduced by Alpher~\cite{Alpher02}, and subsequently developed by
   Alpher and Fotheringham-Smythe~\cite{Alpher03}, and Alpher \etal~\cite{Alpher04}.''

This is incorrect: ``... subsequently developed by Alpher \etal~\cite{Alpher03} ...'' because reference~\cite{Alpher03} has just two authors.


\subsection{Color}

Please refer to the author guidelines on the \confName\ \confYear\ web page for a discussion of the use of color in your document.

If you use color in your plots, please keep in mind that a significant subset of reviewers and readers may have a color vision deficiency; red-green blindness is the most frequent kind. Hence avoid relying only on color as the discriminative feature in plots (such as red \vs green lines), but add a second discriminative feature to ease disambiguation. 


\section{Final copy}

You must include your signed publishing aggreement (digital signatures are encouraged) when you submit your finished paper. We MUST have this form before your paper can be published in the proceedings. The digital form is available via the conference homepage. 


\section*{Acknowledgments}

Are acknowledgments OK?  Yes, but leave them for the final copy.

%%%%%%%%% REFERENCES
{\small
\bibliographystyle{ieee_fullname}
\bibliography{egbib}
}

\end{document}
